% Discussion section

\subsection{Interpretation of findings}

% [Interpretation and context of the results]

This study provides, to our knowledge, the first large-scale automated evaluation of reporting quality in Mendelian randomization (MR) studies using a large language model (LLM) aligned with the STROBE-MR guideline. By analysing over 12,000 MR articles published between 2007 and 2025, several important patterns emerge regarding both the strengths and weaknesses in current reporting practices, as well as the potential role of LLMs in monitoring reporting transparency.

Evaluation of the LLM's performance showed a high accuracy of 89.0\%, demonstrating the feasibility of using LLMs for automated reporting quality assessment. However, the accuracy decreased for items in the Methods and Results sections. This can be explained by the fact that these sections contain more technical details, the LLM has certain knowledge gaps, and the prompt did not provide enough information for the LLM to make correct decisions. For example, item 4c: "For each data source, describe measurement of genetic variants", this item requires articles to report the genotyping details, however, the LLM mixed up the measurement of genetic variants with the selection of genetic variants (p-value, LD clumping, etc.) and thus made wrong decisions. Another example is item 13d: "Report and compare with estimates from non-MR analyses". This item requires articles to include a non-MR analysis as an additional analysis and compare its results with the primary MR analysis. However, the LLM thought articles can just cite other non-MR studies on the same topic and discuss the results, which is not the purpose of this item. These examples highlight the importance of careful prompt design when using LLMs for such tasks. 

Overall, MR studies demonstrate high complicance with some core STROBE-MR items, including research background, key results, findings and meanings of study. 
This suggests that authors are generally effective at communicating the purpose, context, and high-level conclusions of their work. 
However, this constrast with the poor reporting of several methodological and transparency-related items, particularly in Other Information, Methods and Discussion sections. 
According Table ~\ref{tab:report:component}, poor reporting of Other Information happens because authors tend to not provide role of funders and statistical code, which plays an important role in reproducing the results. 
In Methods section, item 4b: paticipants and item 6c: details of MR estimator are rarely reported. 
Item 4 requires the  article to present information on study population while item 6c has extra requirement for two-sample designs, that is to detail the included covariates and whether the same covariates set was used in the samples. 
One explanation for this poor reporting is the dominance of two-sample MR designs (88.6\%)-particularly in recent years. 
These designs rely on secondary data sources, which may lead authors to omit detailed descriptions of participants, data collection, and assumptions, assuming these are already understood or previously reported. 
However, such omissions hinder transparency and reproducibility, especially for readers who are not domain experts. 
In Discussion section, item 17: generalizability and item 16b: Mechanism are most likely missed by authors. 
Table ~\ref{tab:report:component} shows that gene-environment equivalence component of item 16b only has a reporting rate of 10.0\% and generalizability across other exposure periods and levels components of item 17 has reporting rates of 26.2\% and 20.0\%, respectively.  
Gene-environment equivalence is the MR version of consistency assumption. Generalizability across different time periods is also critical because genetic variants can have different effect over the whole life course. Ignoring these items could undermine the validity of findings.

Sensitivity analysis shows a substantial improvement of reporting rates in sections that are considered as poorly reported in main analysis. This can be explained by the fact that many studies do include relevant information but not all that are required by STROBE-MR items. Interstingly, we found that there was no difference in the average number of items reported between studies with and without self-claimed STROBE-MRcompliance. These findings indicate that authors do have different interpretations of STROBE-MR, and the guideline was not used in a correct and consistent way.

These findings suggest that the introduction of the STROBE-MR guideline has not yet translated into uniformly improved reporting practices across the MR literature. Although certain items are well reported, key methodological and transparency elements remain neglected. This may reflect limited awareness of the guideline, lack of enforcement by journals and reviewers, or misunderstanding of certain items. To address these issues, stronger endorsement of STROBE-MR by journals/peer reviewers/funding bodies and more trainings on how to apply STROBE-MR may be necessary.

\subsection{Limitations}

% [Study limitations and potential biases]

There are several limitations in this study. First, the accuracy of LLM-based assessment depends on prompt design and also model selection, we only tried GPT-5-mini, other models may not achieve the same results. 
Second, the reliance on available full texts on PubMed Central may introduce selection bias. 
Third, the study does not directly assess the correctness of reported methods, only their presence or absence. 

% \subsection{Implications}

% [Clinical, policy, or research implications]